\subsection{System Feature 1: Mentor Dashboard} 


\subsubsection{Description}
The mentor shall be able to create and edit events for mentees to view on their
personalized calendars. Mentors can set event visibility so that mentees only
see relevant events they are eligible for. \textit{(Priority: High)} 



\subsubsection{Functional Requirements}

\begin{itemize}[label=$\bullet$]
    \item \textbf{MD-1:} Mentors shall create events with one out of three
    visibility settings (invitation-only, open-to-everyone, or cohort-specific).
    \begin{itemize}
        \item Backward Traceability: This requirement is necessary to implement
        system feature 2, which displays only events that mentee is eligible
        for. This requirement was mentioned in Group 11 RFP Section 1.

        \item Forward Traceability: This requirement can be verified by
        successfully creating an event with a certain visibility.
        Invitation-only events are visible for mentees only if they are invited
        by a mentor, open-to-everyone events are visible to all mentees, and
        cohort-specific events are visible if the mentee belongs to the
        cohort(s) that mentor specified.
    \end{itemize}
    
    \item \textbf{MD-2:} Mentors shall be able to add info about event title,
    location, date, start time, end time, and description.
    \begin{itemize}
        \item Backward Traceability: This requirement was specified by client in
        Client Meeting Notes which occurred on \DTMdate{2026-02-06}.

        \item Forward Traceability: This requirement can be verified by a mentee
        when they click on the event on their personal calendar and are able to
        see info about the event's title, location, date, start time, end time,
        and description.
    \end{itemize}
    
    \item \textbf{MD-3:} Mentors shall be able to edit event title, location,
    date, start time, end time, and description.
    \begin{itemize}
        \item Backward Traceability: Traceability: This requirement was
        specified by client in Group 11 RFP Section 2, Events Management
        Interface.

        \item Forward Traceability: This requirement can be verified by the
        mentor successfully editing details of an event and having the changes
        reflected on mentees' personalized calendars.
    \end{itemize}
    
    \item \textbf{MD-4:} Mentors shall be able to delete events.
    \begin{itemize}
        \item Backward Traceability: This requirement was defined by Group 2 as
        a result of rendering this feature.

        \item Forward Traceability: This requirement can be verified if the
        event gets deleted from the database and all mentees' personal
        calendars.
    \end{itemize}
    
    \item \textbf{MD-5:} Mentors shall have access to the RSVP list of events
    they are hosting.
    \begin{itemize}
        \item Backward Traceability: This requirement was specified in Group 11
        RFP Section 2, Mentor Interface.

        \item Forward Traceability: This requirement can be verified if the
        application displays a list of event attendees when mentor clicks on any
        of their events from the mentor dashboard.
    \end{itemize}
    
    \item \textbf{MD-6:} Mentors shall be able to invite mentees to their
    invitation-only events.
    \begin{itemize}
        \item Backward Traceability: This requirement was specified in the Group
        11 RFP, Glossary of terms section, under Mentor.

        \item Forward Traceability: This requirement can be verified if a mentor
        can select mentees to invite from their dashboard. Personalized calendar
        should display the event as eligible only for mentees that have been
        invited.
    \end{itemize}
\end{itemize}



\subsubsection{Use Cases}

% Use Case 1
\begin{usecase}{UC-1: Create a Event} 
    \textbf{Primary Actor}  & Mentor                                             \\ \hline
    \textbf{Description}    & The mentor creates an event and sets visibility as
                              either invitation-only, open-to-everyone, or
                              cohort-specific. The mentor inputs info regarding
                              event title, date, location, start time, end time,
                              and description.                                   \\ \hline
    \textbf{Trigger}        & Mentor indicates that they want to create an event
                              from their dashboard.                              \\ \hline
    \textbf{Preconditions}  & PRE-1: User is authorized as mentor.               \\ \hline
    \textbf{Postconditions} & POST-1: New event is stored in a database.         \newline 
                              POST-2: New event displayed on calendars of
                              mentees that satisfy visibility settings set forth
                              by mentor.                                         \\ \hline
    \textbf{%
        Normal Flow
        (Open-to-everyone)%
    }                       & 
    \begin{tabenum}
        \item Mentor selects "Create Event" button on their mentor dashboard
        \item Application displays empty event form
        \item Mentor selects open-to-everyone visibility
        \item Mentor inputs details regarding the event's title, location, date, start time, end time, and description
        \item Mentor clicks "Publish Event"
        \item Application saves event to database
        \item Event is displayed as eligible on all mentees' personalized calendars
    \end{tabenum}                                                                \\ \hline
    \textbf{%
        Alternative Flow 1.1
        (Cohort-specific)%
    }                       & 
    \begin{tabenum}
        \item Mentor clicks "Create New Event" button on their mentor dashboard.
        \item Application displays empty event form.
        \item Mentor selects cohort-specific visibility.
        \item Application provides additional input field for cohort (Freshman, sophomore, etc.)
        \item Mentor inputs cohort, event title, location, date, start time, end time, and description.
        \item Mentor clicks "Publish Event"
        \item Application saves event to database
        \item Event is displayed on personal calendars of mentees belonging to specified cohort
    \end{tabenum}                                                                \\ \hline
    \textbf{%
        Alternative Flow 1.2
        (Invitation-only)%
    }                       & 
    \begin{tabenum}
        \item Mentor clicks "Create New Event" button on their mentor dashboard.
        \item Application displays empty event form.
        \item Mentor selects invitation-only visibility
        \item Application displays list of all mentees
        \item Mentor selects mentees they would like to invite from list.
        \item Mentor inputs event title, location, date, start time, end time, and description.
        \item Mentor clicks "Publish Event"
        \item Application saves event to database.
        \item Event is displayed on personal calendars of invited mentees as eligible
    \end{tabenum}                                                                \\ \hline
    \textbf{Exceptions}     & If any one of the fields (visibility, title, location,
                              date, start time, end time, and description) are missing,
                              event creation will be blocked.                    \\ \hline
    \textbf{Priority}       & High                                               \\
\end{usecase}


% Use Case 2
\begin{usecase}{UC-2: Edit Event Details} 
    \textbf{Primary Actor}  & Mentor                                            \\ \hline
    \textbf{Description}    & A mentor can edit any information (event visibility,
                              title, location, date, start time, end time, and
                              description) of an event that they have created.
                              After saving changes, the application should reflect
                              these updates on mentees' personal calendars.     \\ \hline
    \textbf{Trigger}        & Mentor indicates that they want to edit an event. \\ \hline
    \textbf{Preconditions}  & PRE-1: Mentor is the creator of this event.       \\ \hline
    \textbf{Postconditions} & POST-1: Event details are updated in the database.\newline 
                              POST-2: Updated event details should be reflected
                              on relevant mentee's personal calendars.          \\ \hline
    \textbf{%
        Normal Flow%
    }                       & 
    \begin{tabenum}
        \item Mentor selects an event from their dashboard.
        \item Mentor inputs desired changes for the visibility, title, location, date, start time, end time, or description.
        \item Mentor clicks "Save Changes".
        \item Application updates event details in database.
        \item Changes are reflected in relevant mentee's personal calendars.
    \end{tabenum}                                                                \\ \hline
    \textbf{Exceptions}     & If any field (visibility, title, location, date, start time, end time, description) is missing, event editing shall be blocked. \\ \hline
    \textbf{Priority}       & High                                               \\
\end{usecase}


% Use Case 3
\begin{usecase}{UC-3: Delete Event} 
    \textbf{Primary Actor}  & Mentor                                             \\ \hline
    \textbf{Description}    & Mentor can delete an event that they have created.
                              The event should be deleted from the entire
                              application.                                       \\ \hline
    \textbf{Trigger}        & Mentor indicates that they want to delete an event. \\ \hline
    \textbf{Preconditions}  & PRE-1: Mentor is the creator of this event.    \\ \hline
    \textbf{Postconditions} & POST-1: Event is deleted from the database. \newline 
                              POST-2: The event is deleted from the mentees'
                              personalized calendars.                       \\ \hline
    \textbf{%
        Normal Flow%
    }                       & 
    \begin{tabenum}
        \item Mentors select an event on their mentor dashboard.
        \item Mentor clicks on "Delete Event".
        \item Application deletes events from database.
        \item Application removes event from all mentees' personalized calendars.
    \end{tabenum}                                                                \\ \hline
    \textbf{Priority}       & High                                               \\
\end{usecase}


% Use Case 4
\begin{usecase}{UC-4: View RSVP List} 
    \textbf{Primary Actor}  & Mentor                                             \\ \hline
    \textbf{Description}    & Mentors can view attendees on the RSVP list for
                              an event that they have created.                  \\ \hline
    \textbf{Trigger}        & Mentor indicates that they want to view the RSVP
                              list.                                              \\ \hline
    \textbf{Preconditions}  & PRE-1: Mentor is the creator of this event.    \\ \hline
    \textbf{Postconditions} & POST-1: The application displays a list of attendees
                              for an event.                                      \\ \hline
    \textbf{%
        Normal Flow%
    }                       & 
    \begin{tabenum}
        \item Mentors select an event on their mentor dashboard.
        \item Application displays RSVP list.
    \end{tabenum}                                                                \\ \hline
    \textbf{Priority}       & High                                               \\
\end{usecase}


\subsubsection{User Stories}
\begin{itemize}[leftmargin=*, label=$\square$]
    \item \textbf{As a mentor}, I want to create events so that mentees can view
    them on their personalized calendars.
    \begin{itemize}[label=$\circ$, nosep]
        \item \textit{Acceptance:} Mentor can access new event form after
        selecting "Create Event".
        \item \textit{Acceptance:} A mentor can create a new event that is
        displayed on personalized calendars of eligible mentees.
        \item \textit{Acceptance:} The application saves new event details to
        database.
    \end{itemize}

    \item \textbf{As a mentor}, I want to control event visibility so that
    mentees see only the events that they are eligible to attend.
    \begin{itemize}[label=$\circ$, nosep]
        \item \textit{Acceptance:} A mentor can set the visibility of an event
        as one of three settings (invitation-only, open to everyone, or
        cohort-specific).
        \item \textit{Acceptance:} A mentee's personal calendar only displays
        events that they are eligible for.
    \end{itemize}

    \item \textbf{As a mentor}, I want to add event details so that mentees can
    make decision to RSVP or not.
    \begin{itemize}[label=$\circ$, nosep]
        \item \textit{Acceptance:} New event form contains input fields for
        event's title, date, start time, end time, location, and description.
        \item \textit{Acceptance:} A event can not be created if it is missing
        event title, date, start time, end time, location, or description.
    \end{itemize}

    \item \textbf{As a mentor}, I want to edit event details so that I can keep
    mentees updated on any changes.
    \begin{itemize}[label=$\circ$, nosep]
        \item \textit{Acceptance:} A mentor can select an event they created and
        edit the event's title, date, start time, end time, location, or
        description.
        \item \textit{Acceptance:} A mentor can not save edits if any input
        field is missing.
        \item \textit{Acceptance:} After mentor saves changes, event details in
        database should be updated.
        \item \textit{Acceptance:} Editing changes should be reflected on
        mentee's personal calendars.
    \end{itemize}

    \item \textbf{As a mentor}, I want to delete an event so that mentees know
    it is cancelled.
    \begin{itemize}[label=$\circ$, nosep]
        \item \textit{Acceptance:} A mentor can select an event they have
        created and delete the event.
        \item \textit{Acceptance:} Event should be deleted from database.
        \item \textit{Acceptance:} Event should be removed from all mentee's
        personal calendars.
    \end{itemize}

    \item \textbf{As a mentor}, I want to view the RSVP list so that I can
    monitor attendance and plan accordingly.
    \begin{itemize}[label=$\circ$, nosep]
        \item \textit{Acceptance:} A mentor can select an event they have
        created and access RSVP list.
        \item \textit{Acceptance:} RSVP list should contain names of all mentees
        who have RSVP'd.
        \item \textit{Acceptance:} RSVP list should not contain names of mentees
        who have un-RSVP'd.
    \end{itemize}
\end{itemize}
